\chapter{问题一模型与结果：稳健传感器校准}

\section{模型建立}

附件1给出两组同步位移数据。两类传感器测量同一物理量，因此首先考虑仿射校准
\begin{equation}
  y_i=\beta_0+\beta_1x_i+\varepsilon_i ,
\end{equation}
其中$x_i$为数据A，$y_i$为数据B。高阶多项式虽然可能降低训练误差，但会削弱外推稳定性。本文采用Huber损失估计参数：
\begin{equation}
  \min_{\beta_0,\beta_1}\sum_i \rho_\delta(y_i-\beta_0-\beta_1x_i),
  \quad
  \rho_\delta(u)=
  \begin{cases}
    u^2/2, & |u|\leq \delta,\\
    \delta(|u|-\delta/2), & |u|>\delta.
  \end{cases}
\end{equation}
该损失在小残差处保持二次型，在大残差处降为线性惩罚，适合处理传感器偶发误差。Huber估计和Hampel识别均属于稳健统计方法\citep{hampel1974influence}。

\section{模型对照与参数}

为防止只报告单一模型，本文同时比较OLS、Theil--Sen和RANSAC。时间块5折交叉验证结果见表\ref{tab:q1_model_compare}。Huber模型的CV MAE最低且P95AE稳定，故作为主模型。最终参数为
\begin{equation}
  \hat y=-0.5264+0.8817x .
\end{equation}

\begin{table}[htbp]
\centering
\small
\caption{问题一校准模型对照}
\label{tab:q1_model_compare}
\begin{tabular}{lrrrrr}
\toprule
模型 & $\hat\beta_0$ & $\hat\beta_1$ & CV RMSE/mm & CV MAE/mm & CV P95AE/mm \\
\midrule
Huber & -0.526 & 0.8817 & 2.684 & 1.289 & 7.219 \\
RANSAC & -0.450 & 0.8813 & 2.683 & 1.300 & 7.197 \\
OLS & -0.370 & 0.8807 & 2.685 & 1.337 & 7.186 \\
Theil--Sen & -0.380 & 0.8817 & 3.932 & 2.787 & 8.567 \\
\bottomrule
\end{tabular}
\end{table}

\begin{figure}[htbp]
\centering
\includegraphics[width=0.88\textwidth]{q1_calibration_diagnostics.png}
\caption{问题一稳健校准与残差分布}
\label{fig:q1_diag}
\end{figure}

图\ref{fig:q1_diag}显示，校准曲线与两传感器散点的主趋势一致，残差分布以零为中心，仅存在少量尾部误差。因此一次稳健仿射模型足以完成校准。

从模型选择角度看，OLS与Huber的RMSE接近，说明两组传感器关系本身近似线性；但Huber的MAE更低，说明其对大残差点更稳。Theil--Sen在交叉验证中误差明显增大，表明在本数据量和噪声结构下过度依赖秩统计斜率会损失效率。RANSAC虽然表现接近Huber，但其结果依赖内点阈值，参数解释不如Huber连续稳定。因此最终采用Huber作为主模型。

\section{表1.1结果}

将5个待校正值代入主模型，并用bootstrap重采样给出95\%置信区间，结果见表\ref{tab:q1_result}。置信区间宽度均小于0.03 mm，说明校准参数对样本扰动不敏感。

\begin{table}[htbp]
\centering
\small
\caption{待校正位移结果（题目表1.1）}
\label{tab:q1_result}
\begin{tabular}{rrrr}
\toprule
校正前$x$/mm & 校正后$\hat y$/mm & 95\%下界/mm & 95\%上界/mm \\
\midrule
7.132 & 5.762 & 5.749 & 5.774 \\
18.526 & 15.809 & 15.796 & 15.820 \\
84.337 & 73.836 & 73.827 & 73.845 \\
123.554 & 108.415 & 108.407 & 108.422 \\
167.667 & 147.311 & 147.304 & 147.317 \\
\bottomrule
\end{tabular}
\end{table}
